\subsubsection*{Planes in three-dimensional space}

A line allows one independent direction of motion. A plane is produced when the
requirement is enlarged to two independent directions:
\[
\Pi=\{P_0+s\mathbf u+t\mathbf v:s,t\in\mathbb R\}.
\]
This section makes dimension visible through freedom: one parameter generates a
line, while two independent parameters generate a plane. If the directions are
parallel, the apparent two-parameter expression still collapses to one geometric
direction.

The cross product now receives the interpretation prepared in Chapter 2. Since
$\mathbf u\times\mathbf v$ is perpendicular to every linear combination of
$\mathbf u$ and $\mathbf v$, it is forced to be a normal vector of the plane. The
dot product converts this fact into
\[
\mathbf n\cdot(P-P_0)=0.
\]
The parametric and normal forms therefore arise from the same construction: one
lists allowed directions, while the other records that motion in the normal
direction must vanish.

The plane through three non-collinear points, relative positions of planes, the
line of intersection of two planes, angles between planes, and point-to-plane
distance are all developed from this framework. The student should learn to
construct both parametric and normal equations, obtain a normal from a cross
product, determine a plane from three points, classify relative positions through
normal vectors, find the direction of an intersection line, and derive distance
as an orthogonal projection. The intended habit is to reuse earlier ideas rather
than invent a new formula for every new object.
